%% kapa: exam=1 round=28 year=2017 period=2 form=A subject=appraisal_law subject_ko=감정평가관계법규 seq=27 number=27
%% source: Q-Net 감정평가사 기출문제 (cst00309) artlSeq=5199603 file=extracted/1008/5199603/2017년 제28회 감정평가사 1차시험 문제지/2교시 A형.pdf
%% conversion: kordoc 4.24.51 -> tools/gen1.py (auto); review=proofread-sonnet; flags=none
\begin{question}{27}
건축법령상 허가 대상 건축물이라 하더라도 건축신고를 하면 건축허가를 받은 것으로 보는 경우를 모두 고른 것은?
\begin{framedtext}
ㄱ. 연면적이 150 제곱미터이고 2층인 건축물의 대수선 \par
ㄴ. 보를 5개 수선하는 것 \par
ㄷ. 내력벽의 면적을 50 제곱미터 수선하는 것 \par
ㄹ. 소규모 건축물로서 연면적의 합계가 150 제곱미터인 건축물의 신축 \par
ㅁ. 소규모 건축물로서 건축물의 높이를 5 미터 증축하는 건축물의 증축
\end{framedtext}
\begin{choices}
  \item ㄱ, ㄴ, ㄷ
  \item ㄱ, ㄷ, ㄹ
  \item ㄱ, ㄹ, ㅁ
  \item ㄴ, ㄷ, ㄹ
  \item ㄴ, ㄷ, ㄹ, ㅁ
\end{choices}
\end{question}
\answer{1}
